\section{Related Work}

\paragraph{Vietnamese sentiment and emotion resources.}
The three datasets that anchor Vietnamese sentiment evaluation are
UIT-VSFC \citep{vannguyen2018uitvsfc}, a 3-way student-feedback
corpus; UIT-VSMEC \citep{ho2020uitvsmec}, a 7-way emotion corpus over
Facebook comments; and AIVIVN \citep{aivivn2019}, a public 3-way
e-commerce review competition. VLSP-2020 \citep{vlsp2020sentiment}
provided shared-task settings around these resources. Domain
expansion has produced ViHSD \citep{vihsd2021},
ViHOS \citep{vihos2023}, the Vietnamese hate-speech-span benchmark,
ViMMRC for reading comprehension \citep{nguyen2020vimmrc}, and broader
NLU benchmarks (VLUE, ViGLUE)
\citep{vlue2024,viglue2024}. Pragmatic phenomena (implicit sentiment,
sarcasm, irony, mocking, code-switching) have not, to our knowledge,
been jointly labelled in a public Vietnamese corpus; \method{}
directly addresses that gap.

\paragraph{Vietnamese encoders and instruction LMs.}
PhoBERT \citep{phobert2020} remains the standard Vietnamese encoder
and dominates classification leaderboards;
ViSoBERT \citep{visobert2023} adds social-media-domain pretraining.
For instruction-tuned baselines we use Sailor-7B \citep{sailor2024}
(a Southeast-Asian instruction LM) and Vistral-7B \citep{vistral2024}
(a Vietnamese-Mistral derivative), and we use the multilingual encoder
XLM-R \citep{conneau2020xlmr}. We also report GPT-4o-mini
\citep{openai2024gpt4o} in zero-shot and 8-shot prompted modes. The
backbones build on Transformer pretraining
\citep{devlin2019bert,liu2019roberta,jiang2023mistral,touvron2023llama2}.

\paragraph{Multi-task learning.}
Multi-task learning has been used to share statistical strength
across related NLP tasks since \citet{caruana1997mtl} and
\citet{collobert2008nlpscratch}; MT-DNN \citep{liu2019mtdnn} extended
the idea to Transformer encoders, and homoscedastic uncertainty
weighting \citep{kendall2018uncertainty} provides a principled way to
combine heterogeneous loss heads. \method{} uses this stack for a
specific configuration: rare pragmatic targets (sarcasm, mocking) as
the primary objective, with frequent ordinary-sentiment and emotion
labels as auxiliary signal.

\paragraph{Sarcasm, irony, and figurative language.}
Sarcasm and irony detection have been studied extensively in English
\citep{ghosh2017role,ghosh2018sarcasmreasons,joshi2017sarcasmsurvey,vandehey2018irony,barbieri2020tweeteval}
and multimodally \citep{cai2019multimodal}. Figurative-language
benchmarks such as FLUTE \citep{chakrabarty2022figurative} and
IMPLI \citep{stowe2022impli} test whether models can reason about
non-literal meaning. Implicit sentiment specifically has been
isolated and surveyed
\citep{liao2024implicitsurvey,li2021dualgraph,liu2023chatgptimplicit},
including a study that finds prompted LLMs lag dedicated implicit-sentiment
classifiers despite high overall sentiment accuracy. Our work brings
these strands into a Vietnamese, multi-label, multi-task setting and
asks whether dataset-level pragmatic labels plus rationale
distillation are enough to close the gap on a single small encoder.

\paragraph{Code-switching.}
Code-switching detection has been studied at the token-tag level
\citep{solorio2014overview} and benchmarked end-to-end with
GLUECoS \citep{khanuja2020gluecos}. Vietnamese-English social text
mixes are documented in dialect and lexical-normalisation resources
\citep{vimd2024,vlexnorm2024} but, until now, not labelled as a
sentiment-relevant phenomenon.

\paragraph{Rationale-augmented training and chain-of-thought distillation.}
Chain-of-thought prompting \citep{wei2022cot} elicits intermediate
explanations from large LMs; a follow-on line of work distils those
explanations into smaller task models as auxiliary supervision
\citep{hsieh2023distillstep,magister2023teaching,shridhar2023distillingcot,hinton2015distillation}.
GoEmotions \citep{demszky2020goemotions} and SemEval-2018 Affect in
Tweets \citep{mohammad2018semeval} are the major English antecedents
for jointly modelling emotion and sentiment; we adopt this design but
specialise it for Vietnamese pragmatics and use rationales as a
training-only signal so that inference remains cheap.

\paragraph{Calibration and evaluation.}
We use expected calibration error \citep{guo2017calibration,naeini2015ece}
to track over-confidence, paired-bootstrap significance
\citep{koehn2004bootstrap}, and Krippendorff's $\alpha$
\citep{krippendorff2011alpha} together with Cohen's
$\kappa$ \citep{cohen1960kappa} for inter-annotator agreement.
Released artifacts follow data-statement and datasheet practice
\citep{bender2018datastatements,gebru2021datasheets,mitchell2019modelcards,bender2021stochastic}.
